\import{algebra/ring.tex}
\import{set/logical-graph.tex}

\section{Opposite rings}

The opposite ring has the same additive structure and the same unit.  Its
multiplication reverses the order of the two factors.

\begin{definition}\label{opposite_multiplication_op}
    $\oppositeMultiplicationOp:
        \Set\to\Set\to\Set\to\Set
        =\fun{R,x,y:\Set}{\ringMul[R](y,x)}$.
\end{definition}

\begin{definition}\label{opposite_multiplication_graph}
    $\oppositeMultiplicationGraph{R}=
        \Apply{\binaryGraphOnOp}{
            \carrier[R],\carrier[R],
            \Apply{\oppositeMultiplicationOp}{R}}$.
\end{definition}

\begin{proposition}\label{opposite_multiplication_graph_type}
    Let $R$ be a ring.
    Then
        $\oppositeMultiplicationGraph{R}
            \in\funs{
                \carrier[R]\times\carrier[R]}{\carrier[R]}$.
\end{proposition}
\begin{proof}
    Follows by
        \auto{opposite_multiplication_graph,
            opposite_multiplication_op,
            logical_binary_graph_on_function,ring_mul_type}.
\end{proof}

\begin{proposition}\label{opposite_multiplication_graph_apply}
    Let $R$ be a ring and $x,y\in\carrier[R]$.
    Then
        $\apply{\oppositeMultiplicationGraph{R}}{(x,y)}
            =\ringMul[R](y,x)$.
\end{proposition}
\begin{proof}
    Follows by
        \auto{opposite_multiplication_graph,
            opposite_multiplication_op,
            logical_binary_graph_on_apply}.
\end{proof}

\begin{definition}\label{opposite_ring}
    $\OppositeRing{R}=
        \structure[ring]{
            \carrier=\carrier[R],
            \ringAdd=\ringAdd[R],
            \ringZero=\ringZero[R],
            \ringNeg=\ringNeg[R],
            \ringMul=\oppositeMultiplicationGraph{R},
            \ringOne=\ringOne[R]
        }$.
\end{definition}

\begin{proposition}\label{opposite_ring_fields}
    $\carrier[\OppositeRing{R}]=\carrier[R]$ and
    $\ringAdd[\OppositeRing{R}]=\ringAdd[R]$ and
    $\ringZero[\OppositeRing{R}]=\ringZero[R]$ and
    $\ringNeg[\OppositeRing{R}]=\ringNeg[R]$ and
    $\ringMul[\OppositeRing{R}]
        =\oppositeMultiplicationGraph{R}$ and
    $\ringOne[\OppositeRing{R}]=\ringOne[R]$.
\end{proposition}
\begin{proof}
    Follows by
        \auto{opposite_ring,
            ringaggregateprojection_carrier,ringaggregateprojection_ringAdd,
            ringaggregateprojection_ringZero,ringaggregateprojection_ringNeg,
            ringaggregateprojection_ringMul,ringaggregateprojection_ringOne}.
\end{proof}

\begin{proposition}\label{opposite_ring_mul_apply}
    Let $R$ be a ring and $x,y\in\carrier[R]$.
    Then
        $\ringMul[\OppositeRing{R}](x,y)=\ringMul[R](y,x)$.
\end{proposition}
\begin{proof}
    Follows by
        \auto{opposite_ring_fields,
            opposite_multiplication_graph_apply}.
\end{proof}

\begin{theorem}\label{opposite_ring_is_exact}
    Let $R$ be a ring.
    Then $\OppositeRing{R}$ is an exact ring.
\end{theorem}
\begin{proof}
    Let $S=\OppositeRing{R}$.
    $\carrier[S]=\carrier[R]$ and
        $\ringAdd[S]=\ringAdd[R]$ and
        $\ringZero[S]=\ringZero[R]$ and
        $\ringNeg[S]=\ringNeg[R]$ and
        $\ringMul[S]=\oppositeMultiplicationGraph{R}$ and
        $\ringOne[S]=\ringOne[R]$ by
        \auto{opposite_ring_fields}.
    $\ringAdd[S]\in
        \funs{\carrier[S]\times\carrier[S]}{\carrier[S]}$ by
        \obv{ring_add_graph}.
    $\ringZero[S]\in\carrier[S]$ by \obv{ring_zero_type}.
    $\ringNeg[S]\in\funs{\carrier[S]}{\carrier[S]}$ by
        \obv{ring_neg_graph}.
    $\ringMul[S]\in
        \funs{\carrier[S]\times\carrier[S]}{\carrier[S]}$ by
        \auto{opposite_multiplication_graph_type}.
    $\ringOne[S]\in\carrier[S]$ by \obv{ring_one_type}.
    For all $x,y,z\in\carrier[S]$ we have
        $\ringAdd[S](x,\ringAdd[S](y,z))
            =\ringAdd[S](\ringAdd[S](x,y),z)$ by
        \auto{ring_add_associative}.
    For all $x,y\in\carrier[S]$ we have
        $\ringAdd[S](x,y)=\ringAdd[S](y,x)$ by
        \auto{ring_add_commutative}.
    For every $x\in\carrier[S]$ we have
        $\ringAdd[S](\ringZero[S],x)=x$ by
        \auto{ring_add_zero_left}.
    For every $x\in\carrier[S]$ we have
        $\ringAdd[S](\ringNeg[S](x),x)=\ringZero[S]$ by
        \auto{ring_add_inverse_left}.
    Show that for all $x,y,z\in\carrier[S]$ we have
        $\ringMul[S](x,\ringMul[S](y,z))
            =\ringMul[S](\ringMul[S](x,y),z)$.
    \begin{subproof}
        Fix $x,y,z\in\carrier[S]$.
        Hence $x,y,z\in\carrier[R]$.
        $\ringMul[R](z,y),\ringMul[R](y,x)\in\carrier[R]$ by
            \obv{ring_mul_type}.
        $\ringMul[S](y,z)=\ringMul[R](z,y)$ and
            $\ringMul[S](x,y)=\ringMul[R](y,x)$ by
            \auto{opposite_ring_mul_apply}.
        Hence
            $\ringMul[S](x,\ringMul[S](y,z))
                =\ringMul[R](\ringMul[R](z,y),x)$ and
            $\ringMul[S](\ringMul[S](x,y),z)
                =\ringMul[R](z,\ringMul[R](y,x))$ by
            \auto{opposite_ring_mul_apply}.
        Follows by \obv{ring_mul_associative}.
    \end{subproof}
    Show that for every $x\in\carrier[S]$ we have
        $\ringMul[S](\ringOne[S],x)=x$ and
        $\ringMul[S](x,\ringOne[S])=x$.
    \begin{subproof}
        Fix $x\in\carrier[S]$.
        Hence $x\in\carrier[R]$.
        Follows by
            \auto{opposite_ring_mul_apply,
                ring_mul_one_left,ring_mul_one_right}.
    \end{subproof}
    Show that for all $x,y,z\in\carrier[S]$ we have
        $\ringMul[S](x,\ringAdd[S](y,z))
            =\ringAdd[S](\ringMul[S](x,y),\ringMul[S](x,z))$.
    \begin{subproof}
        Fix $x,y,z\in\carrier[S]$.
        Hence $x,y,z\in\carrier[R]$.
        $\ringAdd[R](y,z)\in\carrier[R]$ by
            \obv{ring_add_type}.
        $\ringMul[S](x,\ringAdd[S](y,z))
            =\ringMul[R](\ringAdd[R](y,z),x)$ by
            \auto{opposite_ring_mul_apply}.
        $\ringMul[S](x,y)=\ringMul[R](y,x)$ and
            $\ringMul[S](x,z)=\ringMul[R](z,x)$ by
            \auto{opposite_ring_mul_apply}.
        Follows by \auto{ring_distributive_right}.
    \end{subproof}
    Show that for all $x,y,z\in\carrier[S]$ we have
        $\ringMul[S](\ringAdd[S](x,y),z)
            =\ringAdd[S](\ringMul[S](x,z),\ringMul[S](y,z))$.
    \begin{subproof}
        Fix $x,y,z\in\carrier[S]$.
        Hence $x,y,z\in\carrier[R]$.
        $\ringAdd[R](x,y)\in\carrier[R]$ by
            \obv{ring_add_type}.
        $\ringMul[S](\ringAdd[S](x,y),z)
            =\ringMul[R](z,\ringAdd[R](x,y))$ by
            \auto{opposite_ring_mul_apply}.
        $\ringMul[S](x,z)=\ringMul[R](z,x)$ and
            $\ringMul[S](y,z)=\ringMul[R](z,y)$ by
            \auto{opposite_ring_mul_apply}.
        Follows by \auto{ring_distributive_left}.
    \end{subproof}
    Hence $S$ is a ring by
        \auto{opposite_ring,ringaggregateexact,ring}.
    Follows by \auto{opposite_ring,ringaggregateexact}.
\end{proof}

\begin{theorem}\label{opposite_ring_involutive}
    Let $R$ be a set.
    Suppose $R$ is an exact ring.
    Then $\OppositeRing{\OppositeRing{R}}=R$.
\end{theorem}
\begin{proof}
    Let $S=\OppositeRing{R}$ and $T=\OppositeRing{S}$.
    $S$ is an exact ring by \auto{opposite_ring_is_exact}.
    $T$ is an exact ring by \auto{opposite_ring_is_exact}.
    $\carrier[T]=\carrier[R]$ and
        $\ringAdd[T]=\ringAdd[R]$ and
        $\ringZero[T]=\ringZero[R]$ and
        $\ringNeg[T]=\ringNeg[R]$ and
        $\ringOne[T]=\ringOne[R]$ by
        \auto{opposite_ring_fields}.
    $\ringMul[T]$ and $\ringMul[R]$ are functions and
        $\dom{\ringMul[T]}=\carrier[R]\times\carrier[R]$ and
        $\dom{\ringMul[R]}=\carrier[R]\times\carrier[R]$ by
        \auto{ring_mul_graph,funs_elim}.
    Show that for every $p\in\dom{\ringMul[T]}$ we have
        $\ringMul[T](p)=\ringMul[R](p)$.
    \begin{subproof}
        Fix $p\in\dom{\ringMul[T]}$.
        Hence $p\in\carrier[R]\times\carrier[R]$.
        Take $x,y\in\carrier[R]$ such that $p=(x,y)$ by
            \obv{times_elem_is_tuple}.
        Hence $x,y\in\carrier[S]$ by
            \auto{opposite_ring_fields}.
        $\ringMul[T](x,y)=\ringMul[S](y,x)$ by
            \auto{opposite_ring_mul_apply}.
        $\ringMul[S](y,x)=\ringMul[R](x,y)$ by
            \auto{opposite_ring_mul_apply}.
    \end{subproof}
    Hence $\ringMul[T]=\ringMul[R]$ by \obv{funext_on_domain}.
    Follows by \auto{exact_ring_extensionality}.
\end{proof}
